\documentclass[11pt]{article}
\usepackage[sc]{mathpazo}
\usepackage[T1]{fontenc}
\usepackage[utf8]{inputenc}
\usepackage{amsmath,amssymb}
\usepackage{geometry}
\usepackage{hyperref}
\usepackage{physics}
\usepackage{siunitx}
\usepackage{enumitem}
\geometry{margin=1in}
\hypersetup{colorlinks=true,linkcolor=blue,citecolor=blue,urlcolor=blue}

\title{CatheterSim: Guide to the Brownian--Dynamics Solver}
\author{Repository walkthrough (docs/guide.tex)}
\date{\today}

\begin{document}
\maketitle
\tableofcontents
\newpage

\section{Overview}
This repository simulates active or passive microswimmers navigating catheter-like channels. It extends the \texttt{BD.jl} package with irregular domains, polygonal obstacles, GPU-resident flow fields, and an optional L\'evy reorientation model. Output trajectories are stored in HDF5 and can be visualised or used to train downstream models.

\paragraph{Where to start.} High-level notes live in \texttt{docs/README.md}; the solver internals are summarized in \texttt{docs/solver.md}; driver scripts are listed in \texttt{docs/scripts.md}. This \LaTeX{} document consolidates those pieces into a single guide.

\section{Environment and Dependencies}
\begin{itemize}[leftmargin=*]
  \item Julia environment pinned by \texttt{Project.toml} and \texttt{Manifest.toml}. Tested with Julia~1.7 and CUDA 11.3 (V100 GPUs).
  \item Core packages: \texttt{BD.jl}, \texttt{CUDA.jl}, \texttt{HDF5.jl}, \texttt{StaticArrays.jl}, \texttt{LinearAlgebra}, \texttt{Statistics}, \texttt{Plots.jl} (for \texttt{analysis.jl}).
  \item Activate with \verb|julia --project=.|
    followed by \verb|using Pkg; Pkg.instantiate()|.
\end{itemize}

\section{Geometry and Boundary Conditions}
\subsection{Irregular domain}
The type \texttt{Irregular\{N,T\}} (defined in \texttt{cathetermodule.jl}) subclasses \texttt{BD.AbstractDomain}. It stores Cartesian extents \((x_1,x_2),(y_1,y_2)\) and derived widths \(L_x, L_y\). Only 2D is supported.

\subsection{Obstacles and walls}
Boundary handling is implemented by \texttt{IrregularBC}. For a particle at position $\vb{x}$ and trial displacement $\vb{u}\,\Delta t$:
\begin{enumerate}[leftmargin=*]
  \item \textbf{Tentative step:} $\tilde{\vb{x}} = \vb{x} + \vb{u}\,\Delta t$.
  \item \textbf{No-flux walls:} enforce parallel flat walls in $y$ using \texttt{BD.parallel\_flat\_walls}, yielding a correction $\Delta\vb{x}_{\text{wall}}$.
  \item \textbf{Polygon obstacles:} obstacles are specified as triangles packed in the GPU cache array \texttt{vertices}. For each triangle, the helper \texttt{obstacle} computes signed distances to its edges; if inside, the smallest positive penetration depth sets a reflection displacement $\Delta\vb{x}_{\text{poly}}$.
  \item \textbf{$x$-periodicity:} \texttt{BD.periodicbc} returns both the wrapping shift $\Delta\vb{x}_{\text{per}}$ and image counter increment \(\vb{p}\in\mathbb{Z}^2\).
  \item \textbf{Final position:} $\vb{x}_{n+1}= \tilde{\vb{x}} + \Delta\vb{x}_{\text{wall}} + \Delta\vb{x}_{\text{poly}} + \Delta\vb{x}_{\text{per}}$.
\end{enumerate}

\section{Flow Interpolation and Kernel Override}
\subsection{Structured-grid interpolation}
Given monotone grids $x_{i}$ and $y_{j}$ and GPU fields $U_x,U_y,\Omega$, the helper \texttt{findind} locates bracketing indices with binary search, and \texttt{interpolate\_u} performs bilinear interpolation:
\begin{equation}
  u(x,y) = a_{00} + a_{10} x + a_{01} y + a_{11} xy,
\end{equation}
where the coefficients are assembled from the four surrounding samples $(f_{11},f_{12},f_{21},f_{22})$.

\subsection{Custom Brownian--dynamics kernel}
\texttt{cathetermodule.jl} redefines \texttt{BD.bd\_kernel} for the \texttt{Irregular} domain. For each particle $i$ and batch step:
\begin{align}
  \vb{u}_i &= \underbrace{\vb{U}_{\text{flow}}(\vb{x}_i)}_{\text{interpolated}} + \underbrace{\vb{U}(\vb{x}_i,q_i,t)}_{\text{self-propulsion}} + \underbrace{\vb{\eta}_T}_{\text{Brownian}},\\
  \omega_i &= \Omega_{\text{flow}}(\vb{x}_i) + \Omega(\vb{x}_i,q_i,t) + \eta_R,\\
  q_i^{n+1} &= \mathrm{update\_orientation}(q_i^{n},\omega_i,\Delta t),
\end{align}
then boundary conditions are applied as in Section~\ref{sec:bc}. Data are passed to the supplied callback \texttt{op} (e.g., position writer).

\subsection{Velocity closures}
Provided helpers model different driving conditions:
\begin{itemize}[leftmargin=*]
  \item \texttt{poiseuille\_linear\_velocity\_and\_swim}: $U_x(y)=u_f\bigl(1-4y^2/H^2\bigr) + U_0 \cos q$, $U_y=U_0 \sin q$.
  \item \texttt{linear\_velocity\_swimonly}: self-propulsion only, used when external flow is loaded from file.
  \item \texttt{poiseuille\_angular\_velocity}: $\omega = 4u_f y/H^2$.
  \item \texttt{zero\_angular\_velocity}: convenience when vorticity is precomputed or negligible.
\end{itemize}

\section{L\'evy Reorientation Model}
\label{sec:levy}
\texttt{cathetermodule-levy.jl} augments the kernel with heavy-tailed run times. Each particle stores the next reorientation time in \texttt{ParticleData.data}. When the running clock $t$ exceeds that value:
\begin{enumerate}[leftmargin=*]
  \item Sample a run time $\tau \sim \text{Pareto}(\alpha,\tau_m)$ via
  \begin{equation}
    \tau = \tau_m U^{-1/\alpha}, \qquad U\sim \mathrm{Uniform}(0,1),\quad \tau_m = \frac{(\alpha-1)\tau_R}{\alpha}.
  \end{equation}
  \item Reset orientation with \texttt{random\_orientation} (uniform on $[0,2\pi)$).
  \item Update \texttt{data[i]} \textendash{} the scheduled time for the next reset.
\end{enumerate}
Initialization overloads \texttt{BD.initialize!} to fill positions, orientations, and the \texttt{data} buffer (scalar or per-particle array).

\section{Driver Scripts}
See \texttt{docs/scripts.md} for a concise list. Key entry points in the repository root:
\begin{itemize}[leftmargin=*]
  \item \textbf{\texttt{channel\_interpolate.jl}}: minimal example with two hard-coded triangular obstacles, no imposed flow, swimming-only velocity. Good sanity check for geometry and GPU interpolation.
  \item \textbf{\texttt{channel\_interpolate\_noobstacle.jl}}: baseline Poiseuille channel with \texttt{BD.Box} and \texttt{ChannelBC}. Writes trajectories to \texttt{smooth\_persistent\_U0\dots.h5}.
  \item \textbf{\texttt{channel\_paramsweep.jl}}: sweeps obstacle geometry (\verb|x2,x3,h|), reads flow fields from \texttt{flowdatapath}, writes to \texttt{outpath}. Default grid: $N_x=1001$, $N_y=201$.
  \item \textbf{\texttt{channel\_sweep1--5.jl}}: partitioned parameter sweeps for cluster runs (ranges of \verb|x2| distributed across files).
  \item \textbf{\texttt{channel\_test.jl}}: single geometry for quick debugging; outputs to \texttt{debug/}.
  \item \textbf{\texttt{analysis.jl}}: plots a selected particle track from an HDF5 file and builds a GIF. Update \texttt{file}, \texttt{pid}, \texttt{Nframe}, \texttt{offsets} to match your dataset.
\end{itemize}

\section{Running a Simulation}
\subsection{Setup}
\begin{enumerate}[leftmargin=*]
  \item Generate or copy a flow field text file \verb|flow_{x2}_{x3}_{h}.txt| containing columns: $x$, $y$, $U$, $V$, $\Omega$ on the structured grid.
  \item Edit the chosen driver script to point \texttt{flowdatapath} and \texttt{outpath} to accessible locations on your system.
  \item Adjust computational scale: particle count \texttt{np}, time step \texttt{dt}, total steps \texttt{Nt}, and sampling interval \texttt{sample\_steps}.
\end{enumerate}

\subsection{Execution}
\begin{itemize}[leftmargin=*]
  \item Local test: \verb|julia --project=. channel_interpolate.jl|.
  \item Cluster (SLURM): modify \texttt{jobcatheter.gpu} as needed, then submit with \verb|sbatch jobcatheter.gpu|. The template requests one V100 GPU for 10~hours and loads CUDA~11.3 and Julia~1.7.2.
\end{itemize}

\section{Data Products}
\begin{itemize}[leftmargin=*]
  \item \textbf{HDF5 trajectories}: groups \texttt{/config/\textless frame\textgreater/} store positions (and orientations if requested) every \texttt{sample\_steps}. Filenames encode geometry parameters for downstream identification.
  \item \textbf{GIF previews}: produced by \texttt{analysis.jl} for quick visual inspection.
  \item \textbf{Logs and batch outputs}: \texttt{abp/} and \texttt{levy/} contain SLURM logs and generated data for respective model variants.
\end{itemize}

\section{Repository Layout}
A quick map (mirrors \texttt{docs/structure.md}):
\begin{itemize}[leftmargin=*]
  \item \texttt{cathetermodule.jl}, \texttt{cathetermodule-levy.jl}: solver core.
  \item \texttt{channel\_*.jl}, \texttt{analysis.jl}: drivers and post-processing.
  \item \texttt{abp/}, \texttt{levy/}: experiment batches and outputs.
  \item \texttt{docs/}: Markdown notes and this \LaTeX{} guide.
\end{itemize}

\section{Practical Tips}
\begin{itemize}[leftmargin=*]
  \item GPU memory: the sweeps allocate $N_x\times N_y$ arrays and up to $10^5$ particles; reduce \texttt{np}, \texttt{Nx}, \texttt{Ny}, or \texttt{Nt} for quick smoke tests.
  \item Paths: all shipped scripts use absolute paths specific to the original environment. Update them before running to avoid silent writes elsewhere.
  \item L\'evy runs: ensure \texttt{cache} includes \texttt{alpha} and \texttt{tauR} when using the heavy-tailed model (see \texttt{cathetermodule-levy.jl}).
  \item Debugging geometry: start with \texttt{channel\_test.jl} or \texttt{channel\_interpolate.jl} and plot a few steps to verify obstacle placement.
\end{itemize}

\section{Glossary}
\begin{description}[leftmargin=!,labelwidth=3cm]
  \item[BD] Brownian Dynamics, implemented by \texttt{BD.jl}.
  \item[Irregular] Custom rectangular domain with polygonal obstacles and periodic $x$.
  \item[ParticleData.data] Auxiliary per-particle buffer (used for L\'evy run timers).
  \item[Ux, Uy, Omega] GPU arrays storing flow velocity components and vorticity on a structured grid.
\end{description}

\end{document}
